% Source: Hatcher, "Algebraic Topology", Chapter 0 "Some Underlying Geometric Notions"
% (2002 Cambridge University Press edition), PDF pages 10-29 of the cited source.
% Transcribed from the cited source and organized according to
% leanblueprint conventions (semantic \label, bracketed titles, \uses dependency edges) below.

The aim of this short preliminary chapter is to introduce a few of the most common
geometric concepts and constructions in algebraic topology. The exposition is somewhat
informal, with no theorems or proofs until the last couple pages, and it should be read
in this informal spirit, skipping bits here and there. In fact, this whole chapter could
be skipped now, to be referred back to later for basic definitions.

\textit{To avoid overusing the word `continuous' we adopt the convention that maps
between spaces are always assumed to be continuous unless otherwise stated.}

\section{Homotopy and Homotopy Type}

One of the main ideas of algebraic topology is to consider two spaces to be equivalent
if they have `the same shape' in a sense that is much broader than homeomorphism. To
take an everyday example, the letters of the alphabet can be written either as unions of
finitely many straight and curved line segments, or in thickened forms that are compact
regions in the plane bounded by one or more simple closed curves.



In each case the thin letter is a subspace of the thick letter, and we can continuously
shrink the thick letter to the thin one. A nice way to do this is to decompose a thick
letter, call it $\mathbf{X}$, into line segments connecting each point on the outer boundary
of $\mathbf{X}$ to a unique point of the thin subletter $X$, as indicated in the figure. Then we
can shrink $\mathbf{X}$ to $X$ by sliding each point of $\mathbf{X} - X$ into $X$ along the line segment
that contains it. Points that are already in $X$ do not move.

We can think of this shrinking process as taking place during a time interval
$0 \le t \le 1$, and then it defines a family of functions $f_t : \mathbf{X} \to \mathbf{X}$ parametrized
by $t \in I = [0,1]$, where $f_t(x)$ is the point to which a given point $x \in \mathbf{X}$ has moved
at time $t$. Naturally we would like $f_t(x)$ to depend continuously on both $t$ and $x$,
and this will be true if we have each $x \in \mathbf{X} - X$ move along its line segment at
constant speed so as to reach its image point in $X$ at time $t=1$, while points $x \in X$
are stationary, as remarked earlier.

Examples of this sort lead to the following general definition.

\begin{definition}[Deformation retraction]
\label{def:deformation-retraction}
\group{Homotopy Theory}
\lean{HatcherLib.DeformationRetract}
\leanok
\uses{def:retraction,def:homotopy-rel-a}
A \textbf{deformation retraction} of a space $X$ onto a subspace $A$ is a family of maps
$f_t : X \to X$, $t \in I$, such that $f_0 = \one$ (the identity map), $f_1(X) = A$, and
$f_t|_A = \one$ for all $t$. The family $f_t$ should be continuous in the sense that the
associated map $X \times I \to X$, $(x,t) \mapsto f_t(x)$, is continuous.
\end{definition}

It is easy to produce many more examples similar to the letter examples, with the
deformation retraction $f_t$ obtained by sliding along line segments. The figure on the
left below shows such a deformation retraction of a M\"obius band onto its core circle.



The three figures on the right show deformation retractions in which a disk with two
smaller open subdisks removed shrinks to three different subspaces.

In all these examples the structure that gives rise to the deformation retraction can
be described by means of the following definition.

\begin{definition}[Mapping cylinder]
\label{def:mapping-cylinder}
\group{Homotopy Theory}
\lean{HatcherLib.MappingCylinder}
\leanok
For a map $f : X \to Y$, the \textbf{mapping cylinder} $M_f$ is the quotient space of the
disjoint union $(X \times I) \sqcup Y$ obtained by identifying each $(x,1) \in X \times I$ with
$f(x) \in Y$.
\end{definition}

In the letter examples, the space $X$ is the outer boundary of the thick letter, $Y$ is
the thin letter, and $f : X \to Y$ sends the outer endpoint of each line segment to its
inner endpoint. A similar description applies to the other examples. Then it is a general
fact that a mapping cylinder $M_f$ deformation retracts to the subspace $Y$ by sliding
each point $(x,t)$ along the segment $\{x\} \times I \subset M_f$ to the endpoint $f(x) \in Y$.
Continuity of this deformation retraction is evident in the specific examples above, and
for a general $f : X \to Y$ it can be verified using Proposition A.17 in the Appendix
concerning the interplay between quotient spaces and product spaces.

Not all deformation retractions arise in this simple way from mapping cylinders. For
example, the thick $\mathbf{X}$ deformation retracts to the thin $X$, which in turn deformation
retracts to the point of intersection of its two crossbars. The net result is a deformation
retraction of $\mathbf{X}$ onto a point, during which certain pairs of points follow paths that
merge before reaching their final destination. Later in this section we will describe a
considerably more complicated example, the so-called `house with two rooms'.

\begin{definition}[Homotopy]
\label{def:homotopy}
\group{Homotopy Theory}
\lean{HatcherLib.Homotopy}
\leanok
A \textbf{homotopy} is a family of maps $f_t : X \to Y$, $t \in I$, such that the associated
map $F : X \times I \to Y$ given by $F(x,t) = f_t(x)$ is continuous.
\end{definition}

A deformation retraction $f_t : X \to X$ is a special case of the general notion of a
homotopy (\cref{def:homotopy}).

\begin{definition}[Homotopic maps]
\label{def:homotopic-maps}
\group{Homotopy Theory}
\lean{HatcherLib.Homotopic}
\leanok
\uses{def:homotopy}
Two maps $f_0, f_1 : X \to Y$ are \textbf{homotopic} if there exists a homotopy $f_t$
connecting them, and one writes $f_0 \simeq f_1$.
\end{definition}

\begin{definition}[Retraction]
\label{def:retraction}
\group{Homotopy Theory}
\lean{HatcherLib.IsRetraction}
\leanok
A \textbf{retraction} of $X$ onto a subspace $A$ is a map $r : X \to X$ such that
$r(X) = A$ and $r|_A = \one$. Equivalently, from a more formal viewpoint, a retraction
is a map $r : X \to X$ with $r^2 = r$, since this equation says exactly that $r$ is the
identity on its image. Retractions are the topological analogs of projection operators
in other parts of mathematics.
\end{definition}

In these terms, a deformation retraction of $X$ onto a subspace $A$ is a homotopy
from the identity map of $X$ to a retraction of $X$ onto $A$ (\cref{def:deformation-retraction,def:retraction}).
One could equally well regard a retraction as a map $X \to A$ restricting to the identity
on the subspace $A \subset X$.

Not all retractions come from deformation retractions. For example, a space $X$
always retracts onto any point $x_0 \in X$ via the constant map sending all of $X$ to
$x_0$, but a space that deformation retracts onto a point must be path-connected since
a deformation retraction of $X$ to $x_0$ gives a path joining each $x \in X$ to $x_0$. It is
less trivial to show that there are path-connected spaces that do not deformation retract
onto a point. One would expect this to be the case for the letters `with holes', A, B,
D, O, P, Q, R. In Chapter 1 we will develop techniques to prove this.

\begin{definition}[Homotopy relative to a subspace]
\label{def:homotopy-rel-a}
\group{Homotopy Theory}
\lean{HatcherLib.HomotopyRel}
\leanok
\uses{def:homotopy}
A homotopy $f_t : X \to X$ that gives a deformation retraction of $X$ onto a subspace
$A$ has the property that $f_t|_A = \one$ for all $t$. In general, a homotopy $f_t : X \to Y$
whose restriction to a subspace $A \subset X$ is independent of $t$ is called a
\textbf{homotopy relative to $A$}, or more concisely, a \textbf{homotopy rel $A$}.
\end{definition}

Thus, a deformation retraction of $X$ onto $A$ is a homotopy rel $A$ from the identity
map of $X$ to a retraction of $X$ onto $A$.

If a space $X$ deformation retracts onto a subspace $A$ via $f_t : X \to X$, then if
$r : X \to A$ denotes the resulting retraction and $i : A \to X$ the inclusion, we have
$ri = \one$ and $ir \simeq \one$, the latter homotopy being given by $f_t$. Generalizing this
situation, we make the following definition.

\begin{definition}[Homotopy equivalence and homotopy type]
\label{def:homotopy-equivalence}
\group{Homotopy Theory}
\lean{HatcherLib.HomotopyEquiv}
\leanok
\uses{def:homotopy,def:homotopic-maps}
A map $f : X \to Y$ is called a \textbf{homotopy equivalence} if there is a map
$g : Y \to X$ such that $fg \simeq \one$ and $gf \simeq \one$. The spaces $X$ and $Y$ are said to
be \textbf{homotopy equivalent} or to have the same \textbf{homotopy type}. The notation is
$X \simeq Y$.
\end{definition}

\begin{proposition}[A deformation retraction is a homotopy equivalence]
\label{prop:deformation-retract-homotopy-equivalence}
\group{Homotopy Theory}
\lean{HatcherLib.DeformationRetract.homotopyEquiv}
\leanok
\uses{def:deformation-retraction,def:homotopy-equivalence}
If $X$ deformation retracts onto a subspace $A$, then the inclusion
$i : A \hookrightarrow X$ is a homotopy equivalence; in particular $A \simeq X$.
\end{proposition}

\begin{proof}
\leanok
\uses{def:retraction}
Let $f_t : X \to X$ be the deformation retraction, with terminal retraction
$f_1 = r$ satisfying $r(X) \subset A$ and $r|_A = \one$, and let $i : A \hookrightarrow X$
be the inclusion. Corestricting $r$ to its image gives a map $\bar r : X \to A$.
Because $r|_A = \one$, for every $a \in A$ we have $\bar r\, i(a) = r(a) = a$, so
$\bar r\, i = \one_A$. Conversely $i\, \bar r = r = f_1$, and the family $f_t$ is a
homotopy from $f_0 = \one_X$ to $f_1 = r$, so $i\, \bar r \simeq \one_X$. Thus $i$
and $\bar r$ are inverse homotopy equivalences, and $A \simeq X$.
\end{proof}

It is an easy exercise to check that this is an equivalence relation, in contrast with
the nonsymmetric notion of deformation retraction. For example, the three graphs
consisting of a single circle, a figure-eight-like theta graph, and a square-with-diagonal
graph are all homotopy equivalent since they are deformation retracts of the same
space, as we saw earlier, but none of the three is a deformation retract of any other.

It is true in general that two spaces $X$ and $Y$ are homotopy equivalent if and only
if there exists a third space $Z$ containing both $X$ and $Y$ as deformation retracts.
For the less trivial implication one can in fact take $Z$ to be the mapping cylinder $M_f$
of any homotopy equivalence $f : X \to Y$. We observed previously that $M_f$ deformation
retracts to $Y$, so what needs to be proved is that $M_f$ also deformation retracts to
its other end $X$ if $f$ is a homotopy equivalence. This is shown in
\cref{cor:homotopy-equivalence-iff-deformation-retract-mapping-cylinder}.

\begin{definition}[Contractible space]
\label{def:contractible}
\group{Homotopy Theory}
\lean{HatcherLib.IsContractible}
\leanok
\uses{def:homotopy-equivalence}
A space having the homotopy type of a point is called \textbf{contractible}. This
amounts to requiring that the identity map of the space be \textbf{nullhomotopic}, that is,
homotopic to a constant map. In general, this is slightly weaker than saying the space
deformation retracts to a point; see the exercises at the end of the chapter for an
example distinguishing these two notions.
\end{definition}

\begin{example}[The house with two rooms is contractible]
\label{ex:house-with-two-rooms-contractible}
\group{Homotopy Theory}
\uses{def:contractible,def:mapping-cylinder}
Let us describe now an example of a 2-dimensional subspace of $\R^3$, known as
the \textit{house with two rooms}, which is contractible but not in any obvious way.



To build this space, start with a box divided into two chambers by a horizontal
rectangle, where by a `rectangle' we mean not just the four edges of a rectangle but
also its interior. Access to the two chambers from outside the box is provided by two
vertical tunnels. The upper tunnel is made by punching out a square from the top of the
box and another square directly below it from the middle horizontal rectangle, then
inserting four vertical rectangles, the walls of the tunnel. This tunnel allows entry to
the lower chamber from outside the box. The lower tunnel is formed in similar fashion,
providing entry to the upper chamber. Finally, two vertical rectangles are inserted to
form `support walls' for the two tunnels. The resulting space $X$ thus consists of three
horizontal pieces homeomorphic to annuli plus all the vertical rectangles that form the
walls of the two chambers.

To see that $X$ is contractible, consider a closed $\varepsilon$-neighborhood $N(X)$ of $X$.
This clearly deformation retracts onto $X$ if $\varepsilon$ is sufficiently small. In fact, $N(X)$
is the mapping cylinder of a map from the boundary surface of $N(X)$ to $X$. Less
obvious is the fact that $N(X)$ is homeomorphic to $D^3$, the unit ball in $\R^3$. To see
this, imagine forming $N(X)$ from a ball of clay by pushing a finger into the ball to
create the upper tunnel, then gradually hollowing out the lower chamber, and similarly
pushing a finger in to create the lower tunnel and hollowing out the upper chamber.
Mathematically, this process gives a family of embeddings $h_t : D^3 \to \R^3$ starting with
the usual inclusion $D^3 \hookrightarrow \R^3$ and ending with a homeomorphism onto $N(X)$.

Thus we have $X \simeq N(X) = D^3 \simeq \text{point}$, so $X$ is contractible since homotopy
equivalence is an equivalence relation. In fact, $X$ deformation retracts to a point. For
if $f_t$ is a deformation retraction of the ball $N(X)$ to a point $x_0 \in X$ and if
$r : N(X) \to X$ is a retraction, for example the end result of a deformation retraction
of $N(X)$ to $X$, then the restriction of the composition $r \circ f_t$ to $X$ is a deformation
retraction of $X$ to $x_0$. However, it is quite a challenging exercise to see exactly what
this deformation retraction looks like.
\end{example}

\section{Cell Complexes}
\label{sec:cell-complexes}

A familiar way of constructing the torus $S^1 \times S^1$ is by identifying opposite sides
of a square. More generally, an orientable surface $M_g$ of genus $g$ can be constructed
from a polygon with $4g$ sides by identifying pairs of edges, as shown in the figure in
the first three cases $g = 1,2,3$. The $4g$ edges of the polygon become a union of $2g$
circles in the surface, all intersecting in a single point. The interior of the polygon can
be thought of as an open disk, or a 2-cell, attached to the union of the $2g$ circles. One
can also regard the union of the circles as being obtained from their common point of
intersection, by attaching $2g$ open arcs, or 1-cells. Thus the surface can be built up
in stages: start with a point, attach 1-cells to this point, then attach a 2-cell.



A natural generalization of this is to construct a space by the following procedure.

\begin{definition}[CW complex]
\label{def:cell-complex}
\group{CW Complexes}
\lean{HatcherLib.IsCWComplex}
\leanok
\begin{enumerate}
\item[(1)] Start with a discrete set $X^0$, whose points are regarded as 0-cells.
\item[(2)] Inductively, form the $n$-skeleton $X^n$ from $X^{n-1}$ by attaching $n$-cells
$e_\alpha^n$ via maps $\varphi_\alpha : S^{n-1} \to X^{n-1}$. This means that $X^n$ is the quotient
space of the disjoint union $X^{n-1} \sqcup_\alpha D_\alpha^n$ of $X^{n-1}$ with a collection of $n$-disks
$D_\alpha^n$ under the identifications $x \sim \varphi_\alpha(x)$ for $x \in \partial D_\alpha^n$. Thus as a set,
$X^n = X^{n-1} \sqcup_\alpha e_\alpha^n$ where each $e_\alpha^n$ is an open $n$-disk.
\item[(3)] One can either stop this inductive process at a finite stage, setting $X = X^n$
for some $n < \infty$, or one can continue indefinitely, setting $X = \bigcup_n X^n$. In the
latter case $X$ is given the weak topology: a set $A \subset X$ is open (or closed) iff
$A \cap X^n$ is open (or closed) in $X^n$ for each $n$.
\end{enumerate}
A space $X$ constructed in this way is called a \textbf{cell complex} or \textbf{CW complex}.
The explanation of the letters `CW' is given in the Appendix, where a number of basic
topological properties of cell complexes are proved.
\end{definition}

\begin{definition}[Dimension of a CW complex]
\label{def:dimension-cw-complex}
\group{CW Complexes}
\uses{def:cell-complex}
If $X = X^n$ for some $n$, then $X$ is said to be finite-dimensional, and the smallest
such $n$ is the \textbf{dimension of $X$}, the maximum dimension of cells of $X$.
\end{definition}

\begin{example}[Graphs as 1-dimensional CW complexes]
\label{ex:cw-graph}
\dcref{ch0:0.1}
\group{CW Complexes}
\mathlibok
\uses{def:cell-complex}
A 1-dimensional cell complex $X = X^1$ is what is called a \textbf{graph} in algebraic
topology. It consists of vertices (the 0-cells) to which edges (the 1-cells) are attached.
The two ends of an edge can be attached to the same vertex.
\end{example}

\begin{example}[Euler characteristic of the house with two rooms]
\label{ex:house-two-rooms-euler-characteristic}
\dcref{ch0:0.2}
\group{CW Complexes}
\uses{ex:house-with-two-rooms-contractible,def:cell-complex,thm:euler-characteristic-homology}
The house with two rooms, pictured earlier (\cref{ex:house-with-two-rooms-contractible}),
has a visually obvious 2-dimensional cell complex structure. The 0-cells are the vertices
where three or more of the depicted edges meet, and the 1-cells are the interiors of the
edges connecting these vertices. This gives the 1-skeleton $X^1$, and the 2-cells are the
components of the remainder of the space, $X - X^1$. If one counts up, one finds there
are 29 0-cells, 51 1-cells, and 23 2-cells, with the alternating sum $29 - 51 + 23$ equal
to 1. This is the \textbf{Euler characteristic}, which for a cell complex with finitely many
cells is defined to be the number of even-dimensional cells minus the number of
odd-dimensional cells. As we shall show in \cref{thm:euler-characteristic-homology}, the Euler characteristic of a cell complex depends only on its homotopy
type, so the fact that the house with two rooms has the homotopy type of a point implies
that its Euler characteristic must be 1, no matter how it is represented as a cell complex.
\end{example}

\begin{example}[Minimal CW structure on the sphere $S^n$]
\label{ex:sphere-two-cells}
\dcref{ch0:0.3}
\group{CW Complexes}
\uses{def:cell-complex}
The sphere $S^n$ has the structure of a cell complex with just two cells, $e^0$ and
$e^n$, the $n$-cell being attached by the constant map $S^{n-1} \to e^0$. This is equivalent
to regarding $S^n$ as the quotient space $D^n / \partial D^n$.
\end{example}

\begin{example}[CW structure on real projective space $\RP^n$]
\label{ex:real-projective-space-cw}
\dcref{ch0:0.4}
\group{CW Complexes}
\uses{def:cell-complex,ex:sphere-two-cells}
Real projective $n$-space $\RP^n$ is defined to be the space of all lines through the
origin in $\R^{n+1}$. Each such line is determined by a nonzero vector in $\R^{n+1}$, unique
up to scalar multiplication, and $\RP^n$ is topologized as the quotient space of
$\R^{n+1} - \{0\}$ under the equivalence relation $v \sim \lambda v$ for scalars $\lambda \ne 0$. We can
restrict to vectors of length 1, so $\RP^n$ is also the quotient space $S^n / (v \sim -v)$, the
sphere with antipodal points identified. This is equivalent to saying that $\RP^n$ is the
quotient space of a hemisphere $D^n$ with antipodal points of $\partial D^n$ identified. Since
$\partial D^n$ with antipodal points identified is just $\RP^{n-1}$, we see that $\RP^n$ is obtained
from $\RP^{n-1}$ by attaching an $n$-cell, with the quotient projection $S^{n-1} \to \RP^{n-1}$ as
the attaching map. It follows by induction on $n$ that $\RP^n$ has a cell complex structure
$e^0 \cup e^1 \cup \cdots \cup e^n$ with one cell $e^i$ in each dimension $i \le n$.
\end{example}

\begin{example}[CW structure on $\RP^\infty$]
\label{ex:infinite-real-projective-space}
\dcref{ch0:0.5}
\group{CW Complexes}
\uses{ex:real-projective-space-cw}
Since $\RP^n$ is obtained from $\RP^{n-1}$ by attaching an $n$-cell, the infinite union
$\RP^\infty = \bigcup_n \RP^n$ becomes a cell complex with one cell in each dimension. We can
view $\RP^\infty$ as the space of lines through the origin in $\R^\infty = \bigcup_n \R^n$.
\end{example}

\begin{example}[CW structure on complex projective space $\CP^n$]
\label{ex:complex-projective-space-cw}
\dcref{ch0:0.6}
\group{CW Complexes}
\uses{def:cell-complex}
Complex projective $n$-space $\CP^n$ is the space of complex lines through the origin
in $\C^{n+1}$, that is, 1-dimensional vector subspaces of $\C^{n+1}$. As in the case of $\RP^n$,
each line is determined by a nonzero vector in $\C^{n+1}$, unique up to scalar
multiplication, and $\CP^n$ is topologized as the quotient space of $\C^{n+1} - \{0\}$ under the
equivalence relation $v \sim \lambda v$ for $\lambda \ne 0$. Equivalently, this is the quotient of the
unit sphere $S^{2n+1} \subset \C^{n+1}$ with $v \sim \lambda v$ for $|\lambda| = 1$. It is also possible to obtain
$\CP^n$ as a quotient space of the disk $D^{2n}$ under the identifications $v \sim \lambda v$ for
$v \in \partial D^{2n}$, in the following way. The vectors in $S^{2n+1} \subset \C^{n+1}$ with last
coordinate real and nonnegative are precisely the vectors of the form
$(w, \sqrt{1-|w|^2}) \in \C^n \times \C$ with $|w| \le 1$. Such vectors form the graph of the
function $w \mapsto \sqrt{1-|w|^2}$. This is a disk $D_+^{2n}$ bounded by the sphere
$S^{2n-1} \subset S^{2n+1}$ consisting of vectors $(w,0) \in \C^n \times \C$ with $|w| = 1$. Each vector in
$S^{2n+1}$ is equivalent under the identifications $v \sim \lambda v$ to a vector in $D_+^{2n}$, and the
latter vector is unique if its last coordinate is nonzero. If the last coordinate is zero,
we have just the identifications $v \sim \lambda v$ for $v \in S^{2n-1}$.

From this description of $\CP^n$ as the quotient of $D_+^{2n}$ under the identifications
$v \sim \lambda v$ for $v \in S^{2n-1}$ it follows that $\CP^n$ is obtained from $\CP^{n-1}$ by attaching a
cell $e^{2n}$ via the quotient map $S^{2n-1} \to \CP^{n-1}$. So by induction on $n$ we obtain a
cell structure $\CP^n = e^0 \cup e^2 \cup \cdots \cup e^{2n}$ with cells only in even dimensions.
Similarly, $\CP^\infty$ has a cell structure with one cell in each even dimension.
\end{example}

After these examples we return now to general theory.

\begin{definition}[Characteristic map of a cell]
\label{def:characteristic-map}
\group{CW Complexes}
\uses{def:cell-complex}
Each cell $e_\alpha^n$ in a cell complex $X$ has a \textbf{characteristic map}
$\Phi_\alpha : D_\alpha^n \to X$ which extends the attaching map $\varphi_\alpha$ and is a homeomorphism
from the interior of $D_\alpha^n$ onto $e_\alpha^n$. Namely, we can take $\Phi_\alpha$ to be the
composition $D_\alpha^n \to X^{n-1} \sqcup_\alpha D_\alpha^n \to X^n \hookrightarrow X$ where the middle map is the
quotient map defining $X^n$.
\end{definition}

For example, in the canonical cell structure on $S^n$ described in
\cref{ex:sphere-two-cells}, a characteristic map for the $n$-cell is the quotient map
$D^n \to S^n$ collapsing $\partial D^n$ to a point. For $\RP^n$ a characteristic map for the cell
$e^i$ is the quotient map $D^i \to \RP^i \subset \RP^n$ identifying antipodal points of
$\partial D^i$, and similarly for $\CP^n$.

\begin{definition}[Subcomplex]
\label{def:subcomplex}
\group{CW Complexes}
\mathlibok
\uses{def:cell-complex}
A \textbf{subcomplex} of a cell complex $X$ is a closed subspace $A \subset X$ that is a
union of cells of $X$. Since $A$ is closed, the characteristic map of each cell in $A$ has
image contained in $A$, and in particular the image of the attaching map of each cell
in $A$ is contained in $A$, so $A$ is a cell complex in its own right.
\end{definition}

\begin{definition}[CW pair]
\label{def:cw-pair}
\group{CW Complexes}
\mathlibok
\uses{def:subcomplex}
A pair $(X,A)$ consisting of a cell complex $X$ and a subcomplex $A$ will be
called a \textbf{CW pair}.
\end{definition}

For example, each skeleton $X^n$ of a cell complex $X$ is a subcomplex. Particular
cases of this are the subcomplexes $\RP^k \subset \RP^n$ and $\CP^k \subset \CP^n$ for $k \le n$. These
are in fact the only subcomplexes of $\RP^n$ and $\CP^n$.

There are natural inclusions $S^0 \subset S^1 \subset \cdots \subset S^n$, but these subspheres
are not subcomplexes of $S^n$ in its usual cell structure with just two cells. However, we
can give $S^n$ a different cell structure in which each of the subspheres $S^k$ is a
subcomplex, by regarding each $S^k$ as being obtained inductively from the equatorial
$S^{k-1}$ by attaching two $k$-cells, the components of $S^k - S^{k-1}$. The infinite-dimensional
sphere $S^\infty = \bigcup_n S^n$ then becomes a cell complex as well. Note that the two-to-one
quotient map $S^\infty \to \RP^\infty$ that identifies antipodal points of $S^\infty$ identifies the two
$n$-cells of $S^\infty$ to the single $n$-cell of $\RP^\infty$.

\begin{remark}[Closure of a cell need not be a subcomplex]
\label{rem:closure-of-cell-not-always-subcomplex}
\group{CW Complexes}
\uses{def:subcomplex}
In the examples of cell complexes given so far, the closure of each cell is a
subcomplex, and more generally the closure of any collection of cells is a subcomplex.
Most naturally arising cell structures have this property, but it need not hold in
general. For example, if we start with $S^1$ with its minimal cell structure and attach
to this a 2-cell by a map $S^1 \to S^1$ whose image is a nontrivial subarc of $S^1$, then
the closure of the 2-cell is not a subcomplex since it contains only a part of the 1-cell.
\end{remark}

\section{Operations on Spaces}
\label{sec:operations-on-spaces}

Cell complexes have a very nice mixture of rigidity and flexibility, with enough
rigidity to allow many arguments to proceed in a combinatorial cell-by-cell fashion and
enough flexibility to allow many natural constructions to be performed on them. Here
are some of those constructions.

\begin{definition}[Product CW structure]
\label{def:product-cw-structure}
\group{CW Complexes}
\uses{def:cell-complex}
If $X$ and $Y$ are cell complexes, then $X \times Y$ has the structure of a cell complex
with cells the products $e_\alpha^m \times e_\beta^n$ where $e_\alpha^m$ ranges over the cells of $X$ and
$e_\beta^n$ ranges over the cells of $Y$.
\end{definition}

For example, the cell structure on the torus $S^1 \times S^1$ described at the beginning
of \cref{sec:cell-complexes} is obtained in this way from the standard cell structure on
$S^1$. For completely general CW complexes $X$ and $Y$ there is one small complication:
the topology on $X \times Y$ as a cell complex is sometimes finer than the product
topology, with more open sets than the product topology has, though the two
topologies coincide if either $X$ or $Y$ has only finitely many cells, or if both $X$ and
$Y$ have countably many cells. This is explained in the Appendix. In practice this subtle
issue of point-set topology rarely causes problems, however.

\begin{definition}[Quotient CW structure]
\label{def:quotient-cw-structure}
\group{CW Complexes}
\uses{def:cw-pair}
If $(X,A)$ is a CW pair consisting of a cell complex $X$ and a subcomplex $A$, then
the quotient space $X/A$ inherits a natural cell complex structure from $X$. The cells
of $X/A$ are the cells of $X - A$ plus one new 0-cell, the image of $A$ in $X/A$. For a
cell $e_\alpha^n$ of $X - A$ attached by $\varphi_\alpha : S^{n-1} \to X^{n-1}$, the attaching map for the
corresponding cell in $X/A$ is the composition $S^{n-1} \to X^{n-1} \to X^{n-1}/A^{n-1}$.
\end{definition}

For example, if we give $S^{n-1}$ any cell structure and build $D^n$ from $S^{n-1}$ by
attaching an $n$-cell, then the quotient $D^n/S^{n-1}$ is $S^n$ with its usual cell structure.
As another example, take $X$ to be a closed orientable surface with the cell structure
described at the beginning of \cref{sec:cell-complexes}, with a single 2-cell, and let $A$
be the complement of this 2-cell, the 1-skeleton of $X$. Then $X/A$ has a cell structure
consisting of a 0-cell with a 2-cell attached, and there is only one way to attach a cell
to a 0-cell, by the constant map, so $X/A$ is $S^2$.

\begin{definition}[Suspension]
\label{def:suspension}
\group{Homotopy Theory}
For a space $X$, the \textbf{suspension} $SX$ is the quotient of $X \times I$ obtained by
collapsing $X \times \{0\}$ to one point and $X \times \{1\}$ to another point.
\end{definition}

The motivating example is $X = S^n$, when $SX = S^{n+1}$ with the two `suspension
points' at the north and south poles of $S^{n+1}$, the points $(0,\ldots,0,\pm 1)$. One can
regard $SX$ as a double cone



\begin{definition}[Cone]
\label{def:cone}
\group{Homotopy Theory}
The \textbf{cone} on a space $X$ is $CX = (X \times I)/(X \times \{0\})$.
\end{definition}

The suspension $SX$ is the union of two copies of the cone $CX$ on $X$. If $X$ is a
CW complex, so are $SX$ and $CX$ (\cref{def:suspension,def:cone}) as quotients of $X \times I$
with its product cell structure, $I$ being given the standard cell structure of two 0-cells
joined by a 1-cell.

Suspension becomes increasingly important the farther one goes into algebraic
topology, though why this should be so is certainly not evident in advance. One
especially useful property of suspension is that not only spaces but also maps can be
suspended. Namely, a map $f : X \to Y$ suspends to $Sf : SX \to SY$, the quotient map
of $f \times I : X \times I \to Y \times I$.

\begin{definition}[Join of spaces]
\label{def:join}
\group{Homotopy Theory}
\uses{def:cone,def:suspension}
The cone $CX$ is the union of all line segments joining points of $X$ to an external
vertex, and similarly the suspension $SX$ is the union of all line segments joining
points of $X$ to two external vertices. More generally, given $X$ and a second space
$Y$, one can define the space of all line segments joining points in $X$ to points in $Y$.
This is the \textbf{join} $X * Y$, the quotient space of $X \times Y \times I$ under the
identifications $(x,y_1,0) \sim (x,y_2,0)$ and $(x_1,y,1) \sim (x_2,y,1)$. Thus we are
collapsing the subspace $X \times Y \times \{0\}$ to $X$ and $X \times Y \times \{1\}$ to $Y$.
\end{definition}

For example, if $X$ and $Y$ are both closed intervals, then we are collapsing two
opposite faces of a cube onto line segments so that the cube becomes a tetrahedron. In
the general case, $X * Y$ contains copies of $X$ and $Y$ at its two ends, and every
other point $(x,y,t)$ in $X * Y$ is on a unique line segment joining the point
$x \in X \subset X * Y$ to the point $y \in Y \subset X * Y$, the segment obtained by fixing $x$
and $y$ and letting the coordinate $t$ in $(x,y,t)$ vary.



A nice way to write points of $X * Y$ is as formal linear combinations
$t_1 x + t_2 y$ with $0 \le t_i \le 1$ and $t_1 + t_2 = 1$, subject to the rules $0x + 1y = y$
and $1x + 0y = x$ that correspond exactly to the identifications defining $X * Y$. In
much the same way, an iterated join $X_1 * \cdots * X_n$ can be constructed as the space
of formal linear combinations $t_1 x_1 + \cdots + t_n x_n$ with $0 \le t_i \le 1$ and
$t_1 + \cdots + t_n = 1$, with the convention that terms $0 x_i$ can be omitted.

\begin{definition}[Simplex]
\label{def:simplex}
\group{Homotopy Theory}
\uses{def:join}
A very special case that plays a central role in algebraic topology is when each
$X_i$ is just a point. For example, the join of two points is a line segment, the join of
three points is a triangle, and the join of four points is a tetrahedron. In general, the
join of $n$ points is a convex polyhedron of dimension $n-1$ called a \textbf{simplex}.
Concretely, if the $n$ points are the $n$ standard basis vectors for $\R^n$, then their join
is the $(n-1)$-dimensional simplex
\[
\Delta^{n-1} = \{(t_1,\ldots,t_n) \in \R^n \mid t_1 + \cdots + t_n = 1 \text{ and } t_i \ge 0\}.
\]
Another interesting example is when each $X_i$ is $S^0$, two points. If we take the
two points of $X_i$ to be the two unit vectors along the $i$-th coordinate axis in $\R^n$,
then the join $X_1 * \cdots * X_n$ is the union of $2^n$ copies of the simplex $\Delta^{n-1}$, and
radial projection from the origin gives a homeomorphism between $X_1 * \cdots * X_n$
and $S^{n-1}$.
\end{definition}

If $X$ and $Y$ are CW complexes, then there is a natural CW structure on $X * Y$
having the subspaces $X$ and $Y$ as subcomplexes, with the remaining cells being the
product cells of $X \times Y \times (0,1)$. As usual with products, the CW topology on
$X * Y$ may be finer than the quotient of the product topology on $X \times Y \times I$.

\begin{definition}[Wedge sum]
\label{def:wedge-sum}
\group{Homotopy Theory}
Given spaces $X$ and $Y$ with chosen points $x_0 \in X$ and $y_0 \in Y$, the
\textbf{wedge sum} $X \vee Y$ is the quotient of the disjoint union $X \sqcup Y$ obtained by
identifying $x_0$ and $y_0$ to a single point.
\end{definition}

For example, $S^1 \vee S^1$ is homeomorphic to the figure `8', two circles touching
at a point. More generally one could form the wedge sum $\bigvee_\alpha X_\alpha$ of an arbitrary
collection of spaces $X_\alpha$ by starting with the disjoint union $\bigsqcup_\alpha X_\alpha$ and
identifying points $x_\alpha \in X_\alpha$ to a single point. In case the spaces $X_\alpha$ are cell
complexes and the points $x_\alpha$ are 0-cells, then $\bigvee_\alpha X_\alpha$ is a cell complex since it
is obtained from the cell complex $\bigsqcup_\alpha X_\alpha$ by collapsing a subcomplex to a point.

For any cell complex $X$, the quotient $X^n/X^{n-1}$ is a wedge sum of $n$-spheres
$\bigvee_\alpha S_\alpha^n$, with one sphere for each $n$-cell of $X$.

\begin{definition}[Smash product]
\label{def:smash-product}
\group{Homotopy Theory}
\uses{def:wedge-sum}
Inside a product space $X \times Y$ there are copies of $X$ and $Y$, namely
$X \times \{y_0\}$ and $\{x_0\} \times Y$ for points $x_0 \in X$ and $y_0 \in Y$. These two copies
of $X$ and $Y$ in $X \times Y$ intersect only at the point $(x_0,y_0)$, so their union can be
identified with the wedge sum $X \vee Y$. The \textbf{smash product} $X \wedge Y$ is then
defined to be the quotient $X \times Y / X \vee Y$.
\end{definition}

One can think of $X \wedge Y$ as a reduced version of $X \times Y$ obtained by collapsing
away the parts that are not genuinely a product, the separate factors $X$ and $Y$. The
smash product $X \wedge Y$ is a cell complex if $X$ and $Y$ are cell complexes with $x_0$
and $y_0$ 0-cells, assuming that we give $X \wedge Y$ the cell-complex topology rather than
the product topology in cases when these two topologies differ. For example,
$S^m \wedge S^n$ has a cell structure with just two cells, of dimensions 0 and $m+n$, hence
$S^m \wedge S^n = S^{m+n}$. In particular, when $m = n = 1$ we see that collapsing longitude
and meridian circles of a torus to a point produces a 2-sphere.

\section{Two Criteria for Homotopy Equivalence}
\label{sec:two-criteria-homotopy-equivalence}

Earlier in this chapter the main tool we used for constructing homotopy equivalences
was the fact that a mapping cylinder deformation retracts onto its `target' end. By
repeated application of this fact one can often produce homotopy equivalences between
rather different-looking spaces. However, this process can be a bit cumbersome in
practice, so it is useful to have other techniques available as well. We will describe two
commonly used methods here. The first involves collapsing certain subspaces to points,
and the second involves varying the way in which the parts of a space are put together.

\subsection*{Collapsing Subspaces}

The operation of collapsing a subspace to a point usually has a drastic effect on
homotopy type, but one might hope that if the subspace being collapsed already has
the homotopy type of a point, then collapsing it to a point might not change the
homotopy type of the whole space. Here is a positive result in this direction:

\begin{proposition}[Collapsing a contractible CW subcomplex]
\label{prop:collapsing-contractible-cw-subcomplex}
\group{CW Complexes}
\uses{def:cw-pair,def:contractible}
If $(X,A)$ is a CW pair consisting of a CW complex $X$ and a contractible
subcomplex $A$, then the quotient map $X \to X/A$ is a homotopy equivalence.
\end{proposition}

\begin{remark}[CW special case of collapsing contractible subcomplex]
\label{rem:collapsing-contractible-cw-subcomplex-special-case}
\group{CW Complexes}
\uses{prop:quotient-by-contractible-hep-homotopy-equiv,prop:cw-pairs-have-hep}
This is the CW special case of \cref{prop:quotient-by-contractible-hep-homotopy-equiv}
(since a CW pair automatically has the homotopy extension property, by
\cref{prop:cw-pairs-have-hep}). Hatcher states it here for early use and defers the
proof until that later, more general proposition.
\end{remark}

\begin{example}[Collapsing edges of a graph]
\label{ex:graphs-collapsing-homotopy-equivalent}
\dcref{ch0:0.7}
\group{Homotopy Theory}
\uses{prop:collapsing-contractible-cw-subcomplex}
The three graphs consisting of a single circle, a theta-shaped graph (two vertices
joined by three edges), and a square with one diagonal are homotopy equivalent since
each is a deformation retract of a disk with two holes, but we can also deduce this
from the collapsing criterion above since collapsing the middle edge of the first and
third graphs produces the second graph.

More generally, suppose $X$ is any graph with finitely many vertices and edges. If
the two endpoints of any edge of $X$ are distinct, we can collapse this edge to a point,
producing a homotopy equivalent graph with one fewer edge. This simplification can
be repeated until all edges of $X$ are loops, and then each component of $X$ is either
an isolated vertex or a wedge sum of circles.

This raises the question of whether two such graphs, having only one vertex in each
component, can be homotopy equivalent if they are not in fact just isomorphic graphs.
Exercise 12 at the end of the chapter reduces the question to the case of connected
graphs. Then the task is to prove that a wedge sum $\bigvee_m S^1$ of $m$ circles is not
homotopy equivalent to $\bigvee_n S^1$ if $m \ne n$. This sort of thing is hard to do directly.
What one would like is some sort of algebraic object associated to spaces, depending
only on their homotopy type, and taking different values for $\bigvee_m S^1$ and $\bigvee_n S^1$
if $m \ne n$. In fact the Euler characteristic does this since $\bigvee_m S^1$ has Euler
characteristic $1-m$. But it is a rather nontrivial theorem that the Euler characteristic
of a space depends only on its homotopy type. A different algebraic invariant that
works equally well for graphs, and whose rigorous development requires less effort
than the Euler characteristic, is the fundamental group of a space, the subject of
Chapter 1.
\end{example}

\begin{example}[$S^2$ with two points identified is homotopy equivalent to $S^1 \vee S^2$]
\label{ex:sphere-two-points-identified}
\dcref{ch0:0.8}
\group{Homotopy Theory}
\uses{prop:collapsing-contractible-cw-subcomplex,def:cw-pair}
Consider the space $X$ obtained from $S^2$ by attaching the two ends of an arc $A$
to two distinct points on the sphere, say the north and south poles. Let $B$ be an arc
in $S^2$ joining the two points where $A$ attaches. Then $X$ can be given a CW complex
structure with the two endpoints of $A$ and $B$ as 0-cells, the interiors of $A$ and $B$ as
1-cells, and the rest of $S^2$ as a 2-cell. Since $A$ and $B$ are contractible, $X/A$ and
$X/B$ are homotopy equivalent to $X$.



The space $X/A$ is the quotient $S^2/S^0$, the sphere with two points identified, and
$X/B$ is $S^1 \vee S^2$. Hence $S^2/S^0$ and $S^1 \vee S^2$ are homotopy equivalent, a fact
which may not be entirely obvious at first glance.
\end{example}

\begin{example}[Torus with meridional disks is a necklace of spheres]
\label{ex:torus-meridional-disks-necklace}
\dcref{ch0:0.9}
\group{Homotopy Theory}
\uses{prop:collapsing-contractible-cw-subcomplex}
Let $X$ be the union of a torus with $n$ meridional disks. To obtain a CW structure
on $X$, choose a longitudinal circle in the torus, intersecting each of the meridional
disks in one point. These intersection points are then the 0-cells, the 1-cells are the
rest of the longitudinal circle and the boundary circles of the meridional disks, and
the 2-cells are the remaining regions of the torus and the interiors of the meridional
disks. Collapsing each meridional disk to a point yields a homotopy equivalent space
$Y$ consisting of $n$ 2-spheres, each tangent to its two neighbors, a `necklace with $n$
beads'.



The third space $Z$ in the figure, a strand of $n$ beads with a string joining its two
ends, collapses to $Y$ by collapsing the string to a point, so this collapse is a homotopy
equivalence. Finally, by collapsing the arc in $Z$ formed by the front halves of the
equators of the $n$ beads, we obtain the fourth space $W$, a wedge sum of $S^1$ with $n$
2-spheres. (One can see why a wedge sum is sometimes called a `bouquet' in the older
literature.)
\end{example}

\begin{example}[Reduced suspension]
\label{ex:reduced-suspension}
\dcref{ch0:0.10}
\group{Homotopy Theory}
\uses{def:suspension,def:smash-product}
Let $X$ be a CW complex and $x_0 \in X$ a 0-cell. Inside the suspension $SX$ we
have the line segment $\{x_0\} \times I$, and collapsing this to a point yields a space
$\Sigma X$ homotopy equivalent to $SX$, called the \textbf{reduced suspension} of $X$. For
example, if we take $X$ to be $S^1 \vee S^1$ with $x_0$ the intersection point of the two
circles, then the ordinary suspension $SX$ is the union of two spheres intersecting
along the arc $\{x_0\} \times I$, so the reduced suspension $\Sigma X$ is $S^2 \vee S^2$, a slightly
simpler space. More generally we have $\Sigma(X \vee Y) = \Sigma X \vee \Sigma Y$ for arbitrary CW
complexes $X$ and $Y$. Another way in which the reduced suspension $\Sigma X$ is slightly
simpler than $SX$ is in its CW structure: in $SX$ there are two 0-cells (the two
suspension points) and an $(n+1)$-cell $e^{n+1} \times (0,1)$ for each $n$-cell $e^n$ of $X$,
whereas in $\Sigma X$ there is a single 0-cell and an $(n+1)$-cell for each $n$-cell of $X$
other than the 0-cell $x_0$.

The reduced suspension $\Sigma X$ is actually the same as the smash product
$X \wedge S^1$ since both spaces are the quotient of $X \times I$ with
$X \times \partial I \cup \{x_0\} \times I$ collapsed to a point.
\end{example}

\subsection*{Attaching Spaces}

Another common way to change a space without changing its homotopy type
involves the idea of continuously varying how its parts are attached together.

\begin{definition}[Attaching one space to another]
\label{def:attaching-space}
\group{CW Complexes}
\lean{HatcherLib.AttachingSpace}
\leanok
We start with a space $X_0$ and another space $X_1$ that we wish to attach to $X_0$
by identifying the points in a subspace $A \subset X_1$ with points of $X_0$. The data
needed to do this is a map $f : A \to X_0$, for then we can form a quotient space of
$X_0 \sqcup X_1$ by identifying each point $a \in A$ with its image $f(a) \in X_0$. Let us
denote this quotient space by $X_0 \sqcup_f X_1$, the space $X_0$ with $X_1$
\textbf{attached along $A$ via $f$}. When $(X_1,A) = (D^n, S^{n-1})$ we have the case of
attaching an $n$-cell to $X_0$ via a map $f : S^{n-1} \to X_0$.
\end{definition}

Mapping cylinders are examples of this construction, since the mapping cylinder
$M_f$ of a map $f : X \to Y$ is the space obtained from $Y$ by attaching $X \times I$ along
$X \times \{1\}$ via $f$.

\begin{definition}[Mapping cone]
\label{def:mapping-cone}
\group{Homotopy Theory}
\uses{def:cone,def:attaching-space}
Closely related to the mapping cylinder $M_f$ is the \textbf{mapping cone}
$C_f = Y \sqcup_f CX$ where $CX$ is the cone $(X \times I)/(X \times \{0\})$ and we attach this
to $Y$ along $X \times \{1\}$ via the identifications $(x,1) \sim f(x)$.
\end{definition}

For example, when $X$ is a sphere $S^{n-1}$ the mapping cone $C_f$ is the space
obtained from $Y$ by attaching an $n$-cell via $f : S^{n-1} \to Y$. A mapping cone $C_f$
can also be viewed as the quotient $M_f/X$ of the mapping cylinder $M_f$ with the
subspace $X = X \times \{0\}$ collapsed to a point.



If one varies an attaching map $f$ by a homotopy $f_t$, one gets a family of spaces
whose shape is undergoing a continuous change, it would seem, and one might expect
these spaces all to have the same homotopy type. This is often the case:

\begin{proposition}[Homotopic attaching maps yield homotopy equivalent spaces]
\label{prop:attaching-homotopic-maps-preview}
\group{CW Complexes}
\uses{def:cw-pair,def:attaching-space}
If $(X_1,A)$ is a CW pair and the two attaching maps $f,g : A \to X_0$ are
homotopic, then $X_0 \sqcup_f X_1 \simeq X_0 \sqcup_g X_1$.
\end{proposition}

\begin{remark}[Sharper version of homotopic attaching maps]
\label{rem:attaching-homotopic-maps-sharper-version}
\group{CW Complexes}
\uses{prop:attaching-maps-homotopic-rel-x0}
A sharper, `rel $X_0$' version of this statement, with a full proof, is given later as
\cref{prop:attaching-maps-homotopic-rel-x0}.
\end{remark}

\begin{example}[Re-deriving $S^2/S^0 \simeq S^1 \vee S^2$ via attaching maps]
\label{ex:sphere-two-points-identified-via-attaching}
\dcref{ch0:0.11}
\group{CW Complexes}
\uses{prop:attaching-homotopic-maps-preview,def:mapping-cone}
Let us rederive the result in \cref{ex:sphere-two-points-identified} that a sphere with
two points identified is homotopy equivalent to $S^1 \vee S^2$. The sphere with two
points identified can be obtained by attaching $S^2$ to $S^1$ by a map that wraps a
closed arc $A$ in $S^2$ around $S^1$, as shown in the figure. Since $A$ is contractible,
this attaching map is homotopic to a constant map, and attaching $S^2$ to $S^1$ via a
constant map of $A$ yields $S^1 \vee S^2$. The result then follows since $(S^2,A)$ is a CW
pair, $S^2$ being obtained from $A$ by attaching a 2-cell.


\end{example}

\begin{example}[The necklace as spheres attached to a circle]
\label{ex:necklace-attaching-spheres-circle}
\dcref{ch0:0.12}
\group{CW Complexes}
\uses{prop:attaching-homotopic-maps-preview}
In similar fashion we can see that the necklace in \cref{ex:torus-meridional-disks-necklace}
is homotopy equivalent to the wedge sum of a circle with $n$ 2-spheres. The necklace
can be obtained from a circle by attaching $n$ 2-spheres along arcs, so the necklace is
homotopy equivalent to the space obtained by attaching $n$ 2-spheres to a circle at
points. Then we can slide these attaching points around the circle until they all
coincide, producing the wedge sum.
\end{example}

\begin{example}[Quotient by a subcomplex as a mapping cone]
\label{ex:quotient-contractible-subcomplex-mapping-cone}
\dcref{ch0:0.13}
\group{Homotopy Theory}
\uses{def:cw-pair,def:mapping-cone}
Here is an application of the earlier fact that collapsing a contractible subcomplex
is a homotopy equivalence (\cref{prop:collapsing-contractible-cw-subcomplex}): if
$(X,A)$ is a CW pair, consisting of a cell complex $X$ and a subcomplex $A$, then
$X/A \simeq X \cup CA$, the mapping cone of the inclusion $A \hookrightarrow X$. For we have
$X/A = (X \cup CA)/CA \simeq X \cup CA$ since $CA$ is a contractible subcomplex of
$X \cup CA$.
\end{example}

\begin{example}[Quotient by a subspace contractible in $X$]
\label{ex:quotient-by-contractible-in-x}
\dcref{ch0:0.14}
\group{Homotopy Theory}
\uses{ex:quotient-contractible-subcomplex-mapping-cone,def:suspension}
If $(X,A)$ is a CW pair and $A$ is contractible in $X$, that is, the inclusion
$A \hookrightarrow X$ is homotopic to a constant map, then $X/A \simeq X \vee SA$. Namely, by
\cref{ex:quotient-contractible-subcomplex-mapping-cone} we have $X/A = X \cup CA$, and
then since $A$ is contractible in $X$, the mapping cone $X \cup CA$ of the inclusion
$A \hookrightarrow X$ is homotopy equivalent to the mapping cone of a constant map, which
is $X \vee SA$. For example, $S^n/S^i \simeq S^n \vee S^{i+1}$ for $i < n$, since $S^i$ is
contractible in $S^n$ if $i < n$. In particular this gives $S^2/S^0 \simeq S^2 \vee S^1$, which is
\cref{ex:sphere-two-points-identified} again.
\end{example}

\section{The Homotopy Extension Property}
\label{sec:homotopy-extension-property}

In this final section of the chapter we will actually prove a few things, including
the two criteria for homotopy equivalence described above. The proofs depend upon a
technical property that arises in many other contexts as well. Consider the following
problem. Suppose one is given a map $f_0 : X \to Y$, and on a subspace $A \subset X$ one
is also given a homotopy $f_t : A \to Y$ of $f_0|_A$ that one would like to extend to a
homotopy $f_t : X \to Y$ of the given $f_0$.

\begin{definition}[Homotopy extension property]
\label{def:homotopy-extension-property}
\group{CW Complexes}
\lean{HatcherLib.HasHEP}
\leanok
\uses{def:homotopy}
If the pair $(X,A)$ is such that the extension problem above can always be solved,
one says that $(X,A)$ has the \textbf{homotopy extension property}. Thus $(X,A)$ has
the homotopy extension property if every pair of maps $X \times \{0\} \to Y$ and
$A \times I \to Y$ that agree on $A \times \{0\}$ can be extended to a map $X \times I \to Y$.
\end{definition}

\begin{proposition}[HEP is equivalent to a retraction property]
\label{prop:hep-iff-retract}
\group{CW Complexes}
\lean{HatcherLib.hasHEP_iff_isRetract}
\leanok
\uses{def:homotopy-extension-property,def:retraction}
Let $A$ be a closed subspace of $X$. Then the pair $(X,A)$ has the homotopy
extension property if and only if $X \times \{0\} \cup A \times I$ is a retract of
$X \times I$. (The forward implication holds for an arbitrary subspace $A$; only
the converse uses closedness. For non-closed $A$ the converse still holds by a
more delicate argument, deferred to the Appendix.)
\end{proposition}

\begin{proof}
\leanok
For one implication, the homotopy extension property for $(X,A)$ implies that the
identity map $X \times \{0\} \cup A \times I \to X \times \{0\} \cup A \times I$ extends to a map
$X \times I \to X \times \{0\} \cup A \times I$, so $X \times \{0\} \cup A \times I$ is a retract of $X \times I$.

The converse is equally easy when $A$ is closed in $X$. Then any two maps
$X \times \{0\} \to Y$ and $A \times I \to Y$ that agree on $A \times \{0\}$ combine to give a map
$X \times \{0\} \cup A \times I \to Y$ which is continuous since it is continuous on the closed
sets $X \times \{0\}$ and $A \times I$. By composing this map with a retraction
$X \times I \to X \times \{0\} \cup A \times I$ we get an extension $X \times I \to Y$, so $(X,A)$ has
the homotopy extension property. The hypothesis that $A$ is closed can be avoided by
a more complicated argument given in the Appendix.

If $X \times \{0\} \cup A \times I$ is a retract of $X \times I$ and $X$ is Hausdorff, then $A$ must
in fact be closed in $X$. For if $r : X \times I \to X \times I$ is a retraction onto
$X \times \{0\} \cup A \times I$, then the image of $r$ is the set of points $z \in X \times I$ with
$r(z) = z$, a closed set if $X$ is Hausdorff, so $X \times \{0\} \cup A \times I$ is closed in
$X \times I$ and hence $A$ is closed in $X$.
\end{proof}

In fact a retraction of $X \times I$ onto $X \times \{0\} \cup A \times I$ is automatically a
deformation retraction, by an explicit formula:

\begin{lemma}[A retract $X \times \{0\} \cup A \times I$ of the cylinder is a deformation retract]
\label{lem:hep-cylinder-deformation-retract}
\group{CW Complexes}
\lean{HatcherLib.hepBaseDeformationRetract, HatcherLib.HasHEP.hepBase_deformationRetract}
\leanok
\uses{def:homotopy-extension-property,def:deformation-retraction,prop:hep-iff-retract}
If $X \times \{0\} \cup A \times I$ is a retract of the cylinder $X \times I$, then it is a
deformation retract of it. In particular, if the pair $(X,A)$ has the homotopy
extension property, then $X \times I$ deformation retracts onto
$X \times \{0\} \cup A \times I$.
\end{lemma}

\begin{proof}
\leanok
Let $r : X \times I \to X \times \{0\} \cup A \times I$ be a retraction, with components
$r(x,s) = (r_1(x,s), r_2(x,s))$. Define, for $u \in I$,
\[
h_u(x,s) = \bigl( r_1(x, us),\ (1-u)\,s + u\,r_2(x,s) \bigr).
\]
This is jointly continuous, and it is a homotopy rel $X \times \{0\} \cup A \times I$
from the identity to $r$: at $u = 0$ it is $(r_1(x,0), s) = (x,s)$ since $r$ fixes
$X \times \{0\}$; at $u = 1$ it is $(r_1(x,s), r_2(x,s)) = r(x,s)$; for $a \in A$ the
retraction fixes both $(a, us)$ and $(a, s)$, so
$h_u(a,s) = (a, (1-u)s + us) = (a,s)$; and for $s = 0$ it gives
$h_u(x,0) = (r_1(x,0), u\,r_2(x,0)) = (x, 0)$. Hence $r$, which lands in
$X \times \{0\} \cup A \times I$ and fixes it, is the terminal stage of a deformation
retraction.

For the second statement, the homotopy extension property furnishes the retraction
$r$ by \cref{prop:hep-iff-retract} (whose forward implication needs no closedness
hypothesis on $A$).
\end{proof}

\begin{example}[A pair without the homotopy extension property]
\label{ex:hep-fails-for-bad-subspace}
\group{CW Complexes}
\uses{prop:hep-iff-retract}
A simple example of a pair $(X,A)$ with $A$ closed for which the homotopy
extension property fails is the pair $(I,A)$ where $A = \{0,1,1/2,1/4,1/3,\ldots\}$. It is
not hard to show that there is no continuous retraction
$I \times I \to I \times \{0\} \cup A \times I$. The breakdown of homotopy extension here can be
attributed to the bad structure of $(X,A)$ near $0$. With nicer local structure the
homotopy extension property does hold, as the next example shows.
\end{example}

\begin{example}[Mapping cylinder neighborhoods give the HEP]
\label{ex:mapping-cylinder-neighborhood-hep}
\dcref{ch0:0.15}
\group{CW Complexes}
\uses{def:mapping-cylinder,prop:hep-iff-retract}
A pair $(X,A)$ has the homotopy extension property if $A$ has a mapping cylinder
neighborhood in $X$, by which we mean a closed neighborhood $N$ containing a
subspace $B$, thought of as the boundary of $N$, with $N - B$ an open neighborhood
of $A$, such that there exists a map $f : B \to A$ and a homeomorphism $h : M_f \to N$
with $h|_{A \cup B} = \text{id}$. Mapping cylinder neighborhoods like this occur fairly often.
For example, the thick letters discussed at the beginning of the chapter provide such
neighborhoods of the thin letters, regarded as subspaces of the plane.

To verify the homotopy extension property, notice first that $I \times I$ retracts onto
$I \times \{0\} \cup \partial I \times I$, hence $B \times I \times I$ retracts onto
$B \times I \times \{0\} \cup B \times \partial I \times I$, and this retraction induces a retraction of
$M_f \times I$ onto $M_f \times \{0\} \cup (A \cup B) \times I$. Thus $(M_f, A \cup B)$ has the
homotopy extension property. Hence so does the homeomorphic pair $(N, A \cup B)$.
Now given a map $X \to Y$ and a homotopy of its restriction to $A$, we can take the
constant homotopy on $X - (N-B)$ and then extend over $N$ by applying the
homotopy extension property for $(N, A \cup B)$ to the given homotopy on $A$ and the
constant homotopy on $B$.


\end{example}

\begin{proposition}[CW pairs have the homotopy extension property]
\label{prop:cw-pairs-have-hep}
\dcref{ch0:0.16}
\group{CW Complexes}
\uses{def:cw-pair,def:homotopy-extension-property,def:deformation-retraction}
If $(X,A)$ is a CW pair, then $X \times \{0\} \cup A \times I$ is a deformation retract of
$X \times I$, hence $(X,A)$ has the homotopy extension property.
\end{proposition}

\begin{proof}
There is a retraction $r : D^n \times I \to D^n \times \{0\} \cup \partial D^n \times I$, for example the
radial projection from the point $(0,2) \in D^n \times \R$. Then setting
$r_t = tr + (1-t)\,\mathrm{id}$ gives a deformation retraction of $D^n \times I$ onto
$D^n \times \{0\} \cup \partial D^n \times I$. This deformation retraction gives rise to a deformation
retraction of $X^n \times I$ onto $X^n \times \{0\} \cup (X^{n-1} \cup A^n) \times I$ since
$X^n \times I$ is obtained from $X^n \times \{0\} \cup (X^{n-1} \cup A^n) \times I$ by attaching copies
of $D^n \times I$ along $D^n \times \{0\} \cup \partial D^n \times I$. If we perform the deformation
retraction of $X^n \times I$ onto $X^n \times \{0\} \cup (X^{n-1} \cup A^n) \times I$ during the
$t$-interval $[1/2^{n+1}, 1/2^n]$, this infinite concatenation of homotopies is a
deformation retraction of $X \times I$ onto $X \times \{0\} \cup A \times I$. There is no problem
with continuity of this deformation retraction at $t = 0$ since it is continuous on
$X^n \times I$, being stationary there during the $t$-interval $[0, 1/2^{n+1}]$, and CW
complexes have the weak topology with respect to their skeleta so a map is continuous
iff its restriction to each skeleton is continuous.
\end{proof}



Now we can prove a generalization of the earlier assertion that collapsing a
contractible subcomplex is a homotopy equivalence.

\begin{proposition}[Quotient by a contractible subspace with the HEP]
\label{prop:quotient-by-contractible-hep-homotopy-equiv}
\dcref{ch0:0.17}
\group{CW Complexes}
\lean{HatcherLib.collapseMk_homotopyEquiv}
\leanok
\uses{def:homotopy-extension-property,def:contractible,def:homotopy-equivalence}
If the pair $(X,A)$ satisfies the homotopy extension property and $A$ is
contractible, then the quotient map $q : X \to X/A$ is a homotopy equivalence.
\end{proposition}

\begin{proof}
\leanok
Let $f_t : X \to X$ be a homotopy extending a contraction of $A$, with $f_0 = \one$.
Since $f_t(A) \subset A$ for all $t$, the composition $qf_t : X \to X/A$ sends $A$ to a point
and hence factors as a composition $X \xrightarrow{q} X/A \to X/A$. Denoting the latter
map by $\bar f_t : X/A \to X/A$, we have $qf_t = \bar f_t q$. When $t=1$ we have $f_1(A)$
equal to a point, the point to which $A$ contracts, so $f_1$ induces a map $g : X/A \to X$
with $gq = f_1$. It follows that $qg = \bar f_1$ since
$qg(\bar x) = q f_1(x) = \bar f_1 q(x) = \bar f_1(\bar x)$. The maps $g$ and $q$ are inverse
homotopy equivalences since $gq = f_1 \simeq f_0 = \one$ via $f_t$ and
$qg = \bar f_1 \simeq \bar f_0 = \one$ via $\bar f_t$.
\end{proof}

Another application of the homotopy extension property, giving a slightly more
refined version of one of our earlier criteria for homotopy equivalence, is the
following. Here the definition of $W \simeq Z \rel Y$ for pairs $(W,Y)$ and $(Z,Y)$ is
that there are maps $\varphi : W \to Z$ and $\psi : Z \to W$ restricting to the identity
on $Y$, such that $\psi\varphi \simeq \one$ and $\varphi\psi \simeq \one$ via homotopies
that restrict to the identity on $Y$ at all times. We first prove the statement for an
arbitrary pair $(X_1, A)$ with $A$ closed satisfying the homotopy extension property
(Hatcher's exercise generalizing his CW-pair statement); the CW case follows once CW
pairs are known to have the homotopy extension property.

\begin{lemma}[Homotopic attaching maps and the HEP]
\label{lem:attach-homotopic-hep-rel}
\group{CW Complexes}
\lean{HatcherLib.attachingSpace_homotopyEquiv_rel, HatcherLib.attachingSpace_homotopyEquiv}
\leanok
\uses{def:attaching-space,def:homotopy-extension-property,lem:hep-cylinder-deformation-retract,def:homotopy-rel-a}
Let $A \subset X_1$ be closed, let the pair $(X_1,A)$ satisfy the homotopy extension
property, and let $f,g : A \to X_0$ be homotopic attaching maps. Then
$X_0 \sqcup_f X_1 \simeq X_0 \sqcup_g X_1 \rel X_0$.
\end{lemma}

\begin{proof}
\leanok
Let $F : A \times I \to X_0$ be a homotopy from $f$ to $g$ and form the attaching
cylinder $W = X_0 \sqcup_F (X_1 \times I)$, the cylinder $X_1 \times I$ attached to
$X_0$ along $A \times I$. For each $t_0 \in \{0,1\}$ there is a canonical map
$\alpha_{t_0} : X_0 \sqcup_{F_{t_0}} X_1 \to W$ sending $X_0$ identically to $X_0$ and
$x_1 \mapsto (x_1, t_0)$; it is well defined since $(a, t_0) \sim F(a,t_0)$ in $W$ for
$a \in A$.

By \cref{lem:hep-cylinder-deformation-retract} the homotopy extension property gives a
deformation retraction $(h_u)$ of $X_1 \times I$ onto $X_1 \times \{0\} \cup A \times I$;
conjugating by the flip $(x,t) \mapsto (x, 1-t)$ gives one onto
$X_1 \times \{1\} \cup A \times I$. Fix $t_0$ and let $(h_u)$ be the corresponding
deformation retraction onto $S = X_1 \times \{t_0\} \cup A \times I$, with terminal
retraction $r$. Since each $h_u$ fixes $A \times I$ pointwise — the locus of the
attaching identification — the maps $\one_{X_0} \sqcup h_u$ descend to a homotopy
$\bar h_u : W \to W$ with $\bar h_0 = \one$, constant on $X_0$ and on the level
$X_1 \times \{t_0\}$.

Define the collapse $\gamma : S \to X_0 \sqcup_{F_{t_0}} X_1$ by $\gamma(a, t) = F(a,t)$
for $a \in A$ and $\gamma(x_1, t_0) = x_1$ otherwise; the two prescriptions agree on
$A \times \{t_0\}$ by the attaching identification, and $\gamma$ is continuous because
$A \times I$ and $X_1 \times \{t_0\}$ are closed in $S$ ($A$ being closed). The map
$\beta : W \to X_0 \sqcup_{F_{t_0}} X_1$ given by the identity on $X_0$ and by
$\gamma \circ r$ on the cylinder is then well defined and continuous, and:
$\beta \alpha_{t_0} = \one$ (for $x_1 \in A$ this is the attaching identification, and
$r$ fixes the level $t_0$); and $\alpha_{t_0} \beta = \bar h_1$, which is homotopic to
$\one_W$ rel the image of $\alpha_{t_0}$ via $\bar h_u$. Thus $\alpha_{t_0}$ exhibits
$X_0 \sqcup_{F_{t_0}} X_1$ as a deformation retract of $W$, rel $X_0$.

Applying this at $t_0 = 0$ and $t_0 = 1$ gives $\beta_0, \beta_1$; the compositions
$\varphi = \beta_1 \alpha_0$ and $\psi = \beta_0 \alpha_1$ restrict to the identity on
$X_0$, and $\psi\varphi = \beta_0(\alpha_1\beta_1)\alpha_0 \simeq \beta_0\alpha_0 = \one$
rel $X_0$ via $\bar h$, and symmetrically $\varphi\psi \simeq \one$ rel $X_0$.
\end{proof}

\begin{proposition}[Homotopic attaching maps give a homotopy equivalence rel $X_0$]
\label{prop:attaching-maps-homotopic-rel-x0}
\dcref{ch0:0.18}
\group{CW Complexes}
\uses{def:cw-pair,def:attaching-space,prop:cw-pairs-have-hep,lem:attach-homotopic-hep-rel}
If $(X_1,A)$ is a CW pair and we have attaching maps $f,g : A \to X_0$ that are
homotopic, then $X_0 \sqcup_f X_1 \simeq X_0 \sqcup_g X_1 \rel X_0$.
\end{proposition}

\begin{proof}
A CW pair $(X_1, A)$ has the homotopy extension property by
\cref{prop:cw-pairs-have-hep}, and the subcomplex $A$ is closed in $X_1$. Hence
\cref{lem:attach-homotopic-hep-rel} applies.
\end{proof}

We finish this chapter with a technical result whose proof will involve several
applications of the homotopy extension property. It uses the following stability of
the property under taking products with the interval, which Hatcher invokes inline.

\begin{lemma}[The homotopy extension property is preserved by products with $I$]
\label{lem:hep-product-interval}
\group{CW Complexes}
\lean{HatcherLib.HasHEPMap.prodMap_id}
\leanok
\uses{def:homotopy-extension-property}
If the pair $(X,A)$ has the homotopy extension property, then so does the pair
$(X \times I, A \times I)$.
\end{lemma}

\begin{proof}
\leanok
Write $C(I,Y)$ for the space of continuous maps $I \to Y$. A homotopy, in a fresh
parameter $u \in I$, of a map defined on $X \times I$ is the same datum as a homotopy
in $u$ of a map defined on $X$ with values in $C(I,Y)$: this is the exponential
correspondence between maps $X \times I \times I \to Y$ and maps $X \times I \to C(I,Y)$.
Under it, the extension problem for the pair $(X \times I, A \times I)$ with target $Y$
becomes the extension problem for $(X,A)$ with target $C(I,Y)$, which is solvable by the
homotopy extension property of $(X,A)$. Evaluating the resulting extension at each
point of $I$ recovers the desired map $X \times I \times I \to Y$. (No hypothesis on
$A$, such as closedness, is needed.)
\end{proof}

\begin{proposition}[A homotopy equivalence restricting to the identity on $A$ is a homotopy equivalence rel $A$]
\label{prop:hep-homotopy-equivalence-rel-a}
\dcref{ch0:0.19}
\group{CW Complexes}
\lean{HatcherLib.hep_homotopy_equiv_rel}
\leanok
\uses{def:homotopy-extension-property,def:homotopy-equivalence,lem:hep-product-interval}
Suppose $(X,A)$ and $(Y,A)$ satisfy the homotopy extension property, and
$f : X \to Y$ is a homotopy equivalence with $f|_A = \one$. Then $f$ is a homotopy
equivalence rel $A$.
\end{proposition}

\begin{proof}
\leanok
Let $g : Y \to X$ be a homotopy inverse for $f$. There will be three steps to the proof.

\textbf{(1)} Construct a homotopy from $g$ to a map $g_1$ with $g_1|_A = \one$. Let
$h_t : X \to X$ be a homotopy from $gf = h_0$ to $\one = h_1$. Since $f|_A = \one$, we can
view $h_t|_A$ as a homotopy from $g|_A$ to $\one$. Then since we assume $(Y,A)$ has the
homotopy extension property, we can extend this homotopy to a homotopy
$g_t : Y \to X$ from $g = g_0$ to a map $g_1$ with $g_1|_A = \one$.

\textbf{(2)} Show $g_1 f \simeq \one \rel A$. A homotopy from $g_1 f$ to $\one$ is given by
\[
k_t =
\begin{cases}
g_{1-2t} f, & 0 \le t \le 1/2 \\
h_{2t-1}, & 1/2 \le t \le 1.
\end{cases}
\]
The two definitions agree when $t = 1/2$. Since $f|_A = \one$ and $g_t = h_t$ on $A$, the
homotopy $k_t|_A$ starts and ends with the identity, and its second half simply retraces
its first half, that is, $k_t = k_{1-t}$ on $A$. Define a `homotopy of homotopies'
$k_{tu} : A \to X$ using the parameter domain $I \times I$ for the pairs $(t,u)$: on the
bottom edge define $k_{t0} = k_t|_A$; below a `V'-shaped curve $k_{tu}$ is independent of
$u$, and above it $k_{tu}$ is independent of $t$ (unambiguous since $k_t = k_{1-t}$ on $A$).
Since $k_0 = \one$ on $A$, we have $k_{tu} = \one$ on the left, right, and top edges of the
square. Since $(X,A)$ has the homotopy extension property, so does
$(X \times I, A \times I)$, so we can extend $k_{tu} : A \to X$ to $k_{tu} : X \to X$ with
$k_{t0} = k_t$. Restricting this $k_{tu}$ to the left, top, and right edges of the
$(t,u)$-square gives a homotopy $g_1 f \simeq \one \rel A$.



\textbf{(3)} Since $g_1|_A = \one = f|_A$, we have $fg_1|_A = \one$, so steps (1) and (2) can be
repeated with the pair $(f,g)$ replaced by $(g_1,f)$. The result is a map
$f_1 : X \to Y$ with $f_1|_A = \one$ and $f_1 g_1 \simeq \one \rel A$. Hence
$f_1 \simeq f_1(g_1 f) = (f_1 g_1) f \simeq f \rel A$. From this we deduce that
$f g_1 \simeq f_1 g_1 \simeq \one \rel A$.
\end{proof}

\begin{corollary}[HEP inclusion homotopy equivalence implies deformation retract]
\label{cor:hep-inclusion-homotopy-equiv-deformation-retract}
\dcref{ch0:0.20}
\group{CW Complexes}
\lean{HatcherLib.hep_inclusion_deformation_retract}
\leanok
\uses{prop:hep-homotopy-equivalence-rel-a}
If $(X,A)$ satisfies the homotopy extension property and the inclusion
$A \hookrightarrow X$ is a homotopy equivalence, then $A$ is a deformation retract of $X$.
\end{corollary}

\begin{proof}
\leanok
Apply \cref{prop:hep-homotopy-equivalence-rel-a} to the inclusion $A \hookrightarrow X$,
viewed as the map of pairs $(A,A) \to (X,A)$ whose restriction to $A$ is the identity.
The proposition upgrades the given homotopy inverse to a retraction $r : X \to A$ with
$r|_A = \one$ and $\iota r \simeq \one_X \rel A$, which is exactly a deformation
retraction of $X$ onto $A$.
\end{proof}

The corollary below needs the homotopy extension property of the pair $(M_f, X)$,
where $X$ sits inside the mapping cylinder as the bottom $X \times \{0\}$. Hatcher
obtains it from the mapping cylinder neighborhood $X \times [0,1/2]$ as in
\cref{ex:mapping-cylinder-neighborhood-hep}; we record it as a separate lemma, proved
by an explicit projection formula.

\begin{lemma}[The pair $(M_f, X)$ has the homotopy extension property]
\label{lem:mapping-cylinder-pair-hep}
\group{CW Complexes}
\lean{HatcherLib.hasHEPMap_mcylInclX}
\leanok
\uses{def:mapping-cylinder,def:homotopy-extension-property}
For any map $f : X \to Y$, the pair $(M_f, X)$ — the mapping cylinder together with
the copy $X \times \{0\}$ of its domain — has the homotopy extension property.
\end{lemma}

\begin{proof}
\leanok
Write points of the cylinder part of $M_f$ as $[x,s]$ with $x \in X$, $s \in I$, so
that $[x,1] = [f(x)] \in Y$ and $X$ is included as the bottom $x \mapsto [x,0]$. Let
$\varphi : M_f \to Z$ be a map and $h : X \times I \to Z$ a homotopy of
$\varphi|_X$, so $h(x,0) = \varphi[x,0]$.

Consider the square $I \times I$ with coordinates $(s,t)$ ($s$ the cylinder
coordinate, $t$ the homotopy time), and project it radially from the external point
$(1,2)$ onto the union of the bottom edge $I \times \{0\}$ and the left edge
$\{0\} \times I$: the ray from $(1,2)$ through $(s,t)$ exits the square through the
left edge at height $\beta(s,t) = (t-2s)/(1-s)$ when $t \ge 2s$, and through the
bottom edge at abscissa $\alpha(s,t) = (2s-t)/(2-t)$ when $t \le 2s$. Both formulas
are continuous on their (closed) regions — on the region $t \ge 2s$ one has
$s \le 1/2$, so the denominator $1-s$ is bounded below by $1/2$, while $2-t \ge 1$
everywhere — and on the seam $t = 2s$ both edges project to the common corner
$(0,0)$.

Define $F : M_f \times I \to Z$ by
\[
F([x,s],t) =
\begin{cases}
\varphi[x, \alpha(s,t)] & t \le 2s,\\
h(x, \beta(s,t)) & t \ge 2s,
\end{cases}
\qquad
F([y],t) = \varphi[y] \quad (y \in Y).
\]
The two prescriptions agree on the seam $t = 2s$, where both read
$\varphi[x,0] = h(x,0)$, so $F$ is continuous on (the image of) each of the two
closed pieces of the cylinder part, hence on their union by the pasting lemma. $F$ is
well defined on the quotient $M_f$: the whole right edge of the square projects to
the corner $(1,0)$ — for $t \le 2$ always $\alpha(1,t) = 1$ — so
$F([x,1],t) = \varphi[x,1] = \varphi[f(x)] = F([f(x)],t)$, matching the constant
prescription on $Y$; since $M_f \times I$ is a quotient of
$((X \times I) \sqcup Y) \times I$ (the interval being locally compact), continuity
descends.

Finally, the bottom edge of the square is fixed by the projection
($\alpha(s,0) = s$), so $F(\cdot,0) = \varphi$, and the left edge is fixed
($\beta(0,t) = t$), so $F([x,0],t) = h(x,t)$: $F$ is the required extension of the
homotopy $h$ of $\varphi$. (The mapping cylinder neighborhood of
\cref{ex:mapping-cylinder-neighborhood-hep} is implicit: the region $t \ge 2s$ where
$F$ reads the homotopy $h$ only meets the collar $s \le 1/2$.)
\end{proof}

\begin{corollary}[Homotopy equivalences are exactly deformation retractions of mapping cylinders]
\label{cor:homotopy-equivalence-iff-deformation-retract-mapping-cylinder}
\dcref{ch0:0.21}
\group{CW Complexes}
\lean{HatcherLib.isHmtpyEquiv_iff_isMcylDeformationRetract, HatcherLib.homotopyEquiv_iff_common_deformationRetract}
\leanok
\uses{cor:hep-inclusion-homotopy-equiv-deformation-retract,def:mapping-cylinder,lem:mapping-cylinder-pair-hep}
A map $f : X \to Y$ is a homotopy equivalence iff $X$ is a deformation retract of the
mapping cylinder $M_f$. Hence, two spaces $X$ and $Y$ are homotopy equivalent iff
there is a third space containing both $X$ and $Y$ as deformation retracts.
\end{corollary}

\begin{proof}
\leanok
Let $i : X \to M_f$ be the natural inclusion and $r : M_f \to Y$ the canonical
retraction of $M_f$ onto $Y$, so that $f = ri$. Since $r$ is a homotopy equivalence,
$f$ is a homotopy equivalence iff $i$ is: if $i$ is one then so is the composite
$f = ri$; conversely, if $f$ is one then, $j : Y \to M_f$ denoting the inclusion,
$jf$ is a homotopy equivalence and $i \simeq jri = jf$ (via the deformation
retraction of $M_f$ onto $Y$), so $i$ is a homotopy equivalence, being homotopic to
one. Here we use that the composition of two homotopy equivalences is a homotopy
equivalence and that a map homotopic to a homotopy equivalence is a homotopy
equivalence.

If $X$ is a deformation retract of $M_f$, the inclusion $i$ is a homotopy
equivalence, hence so is $f$. Conversely, if $f$ is a homotopy equivalence then so is
$i$, and applying \cref{cor:hep-inclusion-homotopy-equiv-deformation-retract} to the
pair $(M_f, X)$ — which satisfies the homotopy extension property by
\cref{lem:mapping-cylinder-pair-hep} — upgrades a homotopy inverse of $i$ to a
deformation retraction of $M_f$ onto $X$.

For the final statement: if $X$ and $Y$ are homotopy equivalent via $f$, the mapping
cylinder $M_f$ contains both $X$ and $Y$ as deformation retracts — $X$ by the above,
$Y$ by the sliding deformation retraction of \cref{def:mapping-cylinder}. Conversely,
if some space $Z$ contains both $X$ and $Y$ as deformation retracts, then both
inclusions $X \to Z$ and $Y \to Z$ are homotopy equivalences, and composing the first
with a homotopy inverse of the second gives a homotopy equivalence $X \to Y$.
\end{proof}

\section*{Exercises}
% Exercises are transcribed for completeness but are not given blueprint labels or
% \uses edges, since they are problems for the reader rather than declarations used
% elsewhere in the development.

\begin{enumerate}

\item Construct an explicit deformation retraction of the torus with one point deleted
onto a graph consisting of two circles intersecting in a point, namely, longitude and
meridian circles of the torus.

\item Construct an explicit deformation retraction of $\R^n - \{0\}$ onto $S^{n-1}$.

\item
\begin{enumerate}
\item[(a)] Show that the composition of homotopy equivalences $X \to Y$ and
$Y \to Z$ is a homotopy equivalence $X \to Z$. Deduce that homotopy equivalence is
an equivalence relation.
\item[(b)] Show that the relation of homotopy among maps $X \to Y$ is an
equivalence relation.
\item[(c)] Show that a map homotopic to a homotopy equivalence is a homotopy
equivalence.
\end{enumerate}

\item A \textbf{deformation retraction in the weak sense} of a space $X$ to a subspace
$A$ is a homotopy $f_t : X \to X$ such that $f_0 = \mathrm{id}$, $f_1(X) \subset A$, and
$f_t(A) \subset A$ for all $t$. Show that if $X$ deformation retracts to $A$ in this weak
sense, then the inclusion $A \hookrightarrow X$ is a homotopy equivalence.

\item Show that if a space $X$ deformation retracts to a point $x \in X$, then for each
neighborhood $U$ of $x$ in $X$ there exists a neighborhood $V \subset U$ of $x$ such
that the inclusion map $V \hookrightarrow U$ is nullhomotopic.

\item
\begin{enumerate}
\item[(a)] Let $X$ be the subspace of $\R^2$ consisting of the horizontal segment
$[0,1] \times \{0\}$ together with all the vertical segments $\{r\} \times [0,1-r]$ for $r$ a
rational number in $[0,1]$. Show that $X$ deformation retracts to any point in the
segment $[0,1] \times \{0\}$, but not to any other point. [See the preceding problem.]



\item[(b)] Let $Y$ be the subspace of $\R^2$ that is the union of an infinite number
of copies of $X$ arranged as in the figure below. Show that $Y$ is contractible but
does not deformation retract onto any point.



\item[(c)] Let $Z$ be the zigzag subspace of $Y$ homeomorphic to $\R$ indicated by
the heavier line. Show there is a deformation retraction in the weak sense (see
Exercise 4) of $Y$ onto $Z$, but no true deformation retraction.
\end{enumerate}

\item Fill in the details in the following construction from \cite{Edwards1999} of a
compact space $Y \subset \R^3$ with the same properties as the space $Y$ in Exercise 6,
that is, $Y$ is contractible but does not deformation retract to any point. To begin, let
$X$ be the union of an infinite sequence of cones on the Cantor set arranged
end-to-end, as in the figure. Next, form the one-point compactification of
$X \times \R$. This embeds in $\R^3$ as a closed disk with curved `fins' attached along
circular arcs, and with the one-point compactification of $X$ as a cross-sectional slice.
The desired space $Y$ is then obtained from this subspace of $\R^3$ by wrapping one
more cone on the Cantor set around the boundary of the disk.



\item For $n > 2$, construct an $n$-room analog of the house with two rooms.

\item Show that a retract of a contractible space is contractible.

\item Show that a space $X$ is contractible iff every map $f : X \to Y$, for arbitrary
$Y$, is nullhomotopic. Similarly, show $X$ is contractible iff every map $f : Y \to X$ is
nullhomotopic.

\item Show that $f : X \to Y$ is a homotopy equivalence if there exist maps
$a, h : Y \to X$ such that $fa \simeq \mathrm{id}$ and $hf \simeq \mathrm{id}$. More generally, show that
$f$ is a homotopy equivalence if $fa$ and $hf$ are homotopy equivalences.

\item Show that a homotopy equivalence $f : X \to Y$ induces a bijection between
the set of path-components of $X$ and the set of path-components of $Y$, and that $f$
restricts to a homotopy equivalence from each path-component of $X$ to the
corresponding path-component of $Y$. Prove also the corresponding statements with
components instead of path-components. Deduce that if the components of a space
$X$ coincide with its path-components, then the same holds for any space $Y$
homotopy equivalent to $X$.

\item Show that any two deformation retractions $r_0^1$ and $r_1^1$ of a space $X$
onto a subspace $A$ can be joined by a continuous family of deformation retractions
$r_s^1$, $0 \le s \le 1$, of $X$ onto $A$, where continuity means that the map
$X \times I \times I \to X$ sending $(x,s,t)$ to $r_s^1(x)$ is continuous.

\item Given positive integers $v$, $e$, and $f$ satisfying $v - e + f = 2$, construct a
cell structure on $S^2$ having $v$ 0-cells, $e$ 1-cells, and $f$ 2-cells.

\item Enumerate all the subcomplexes of $S^\infty$, with the cell structure on
$S^\infty$ that has $S^n$ as its $n$-skeleton.

\item Show that $S^\infty$ is contractible.

\item
\begin{enumerate}
\item[(a)] Show that the mapping cylinder of every map $f : S^1 \to S^1$ is a CW
complex.
\item[(b)] Construct a 2-dimensional CW complex that contains both an annulus
$S^1 \times I$ and a M\"obius band as deformation retracts.
\end{enumerate}

\item Show that $S^1 * S^1 = S^3$, and more generally $S^m * S^n = S^{m+n+1}$.

\item Show that the space obtained from $S^2$ by attaching $n$ 2-cells along any
collection of $n$ circles in $S^2$ is homotopy equivalent to the wedge sum of $n+1$
2-spheres.

\item Show that the subspace $X \subset \R^3$ formed by a Klein bottle intersecting
itself in a circle, as shown in the figure, is homotopy equivalent to
$S^1 \vee S^1 \vee S^2$.



\item If $X$ is a connected Hausdorff space that is a union of a finite number of
2-spheres, any two of which intersect in at most one point, show that $X$ is
homotopy equivalent to a wedge sum of $S^1$'s and $S^2$'s.

\item Let $X$ be a finite graph lying in a half-plane $P \subset \R^3$ and intersecting
the edge of $P$ in a subset of the vertices of $X$. Describe the homotopy type of the
`surface of revolution' obtained by rotating $X$ about the edge line of $P$.

\item Show that a CW complex is contractible if it is the union of two contractible
subcomplexes whose intersection is also contractible.

\item Let $X$ and $Y$ be CW complexes with 0-cells $x_0$ and $y_0$. Show that the
quotient spaces $X*Y/(X*\{y_0\} \cup \{x_0\}*Y)$ and
$S(X \wedge Y)/S(\{x_0\} \wedge \{y_0\})$ are homeomorphic, and deduce that
$X*Y \simeq S(X \wedge Y)$.

\item If $X$ is a CW complex with components $X_\alpha$, show that the suspension
$SX$ is homotopy equivalent to $Y \vee \bigvee_\alpha SX_\alpha$ for some graph $Y$. In the case
that $X$ is a finite graph, show that $SX$ is homotopy equivalent to a wedge sum of
circles and 2-spheres.

\item Use \cref{cor:hep-inclusion-homotopy-equiv-deformation-retract} to show that if
$(X,A)$ has the homotopy extension property, then $X \times I$ deformation retracts to
$X \times \{0\} \cup A \times I$. Deduce from this that
\cref{prop:attaching-maps-homotopic-rel-x0} holds more generally for any pair
$(X_1,A)$ satisfying the homotopy extension property.

\item Given a pair $(X,A)$ and a homotopy equivalence $f : A \to B$, show that the
natural map $X \to B \sqcup_f X$ is a homotopy equivalence if $(X,A)$ satisfies the
homotopy extension property. [Hint: consider $X \cup M_f$ and use the preceding
problem.] An interesting case is when $f$ is a quotient map, hence the map
$X \to B \sqcup_f X$ is the quotient map identifying each set $f^{-1}(b)$ to a point. When
$B$ is a point this gives another proof of
\cref{prop:quotient-by-contractible-hep-homotopy-equiv}.

\item Show that if $(X_1,A)$ satisfies the homotopy extension property, then so does
every pair $(X_0 \sqcup_f X_1, X_0)$ obtained by attaching $X_1$ to a space $X_0$ via a
map $f : A \to X_0$.

\item In case the CW complex $X$ is obtained from a subcomplex $A$ by attaching a
single cell $e^n$, describe exactly what the extension of a homotopy $f_t : A \to Y$ to
$X$ given by the proof of \cref{prop:cw-pairs-have-hep} looks like. That is, for a point
$x \in e^n$, describe the path $f_t(x)$ for the extended $f_t$.

\end{enumerate}
